% ==============================================================================
% PLANTILLA MAESTRA GENERAL - UCNL
% ==============================================================================
% Uso:
% 1. Duplica este archivo y renombra la copia por materia, semana y actividad.
% 2. Actualiza documenttitle, documentsubtitle, documentsubject, coursename,
%    coursecode, universitydepartment y datos de la figura docente.
% 3. Lee la planeacion semanal y NO pegues las instrucciones completas dentro del
%    producto final. Usa la planeacion como lista de cumplimiento.
% 4. Revisa la bibliografia indicada, toma notas, registra citas y redacta con
%    comprension propia.
% 5. Identifica el producto academico solicitado y construyelo dentro del
%    documento siempre que sea posible.
% 6. Entrega en PDF con la nomenclatura indicada en el aula.
%
% Formato institucional recomendado:
% - Fuente tipo Helvetica, equivalente visual a Arial.
% - Interlineado 1.5.
% - Citas autor-anio con natbib: usar \citep{clave} o \citet{clave}.
% - Referencias en bibliografia BibTeX, formato APA 7 en cuerpo y referencias.
% - Caratula, resumen breve, desarrollo, producto solicitado, conclusion y
%   bibliografia.
% - Redaccion academica clara, coherente, cohesionada, formal y sin faltas.
% - La portada puede llevar una marca de agua institucional discreta; ajustar
%   las variables coverwatermark... si se requiere apagarla o cambiar intensidad.
%
% Criterios generales de cumplimiento:
% - Responder exactamente a la actividad solicitada en la planeacion.
% - Identificar conceptos, elementos, criterios o categorias de forma precisa.
% - Analizar, interpretar y tomar postura; no solo copiar fuentes.
% - Organizar la informacion por jerarquia, secuencia y criterios claros.
% - Incluir conclusion personal cuando la actividad lo permita o lo solicite.
% - Usar citas y referencias en APA 7.
% - Evitar copiar las indicaciones de la actividad en el cuerpo final.
%
% Regla para productos visuales:
% - Si la actividad pide tabla, cuadro comparativo, esquema, mapa conceptual,
%   mapa mental, linea de tiempo, diagrama, matriz, lista de verificacion, red
%   conceptual, SQA, organizador grafico o producto similar, realizarlo
%   preferentemente con LaTeX/TikZ para mantener nitidez, evidencia editable y
%   diseno controlado.
% - Si la consigna exige Canva, Miro u otra herramienta externa, insertar una
%   captura nitida y conservar la explicacion dentro del documento.
% - Para tablas extensas, usar pagina limpia, sin encabezados ni pies, landscape,
%   letra pequena y restaurar el estilo al terminar.
%
% Criterios de redaccion personal:
% - No escribir para llenar espacio; escribir para sostener una idea.
% - Convertir informacion en criterio: que significa, como funciona, que problema
%   resuelve, que limite presenta y que consecuencia produce.
% - Estructura mental recomendada: problema -> sistema -> consecuencia -> postura.
% - Usar pensamiento de diagnostico: entrada -> proceso -> salida.
% - Cada parrafo debe tener funcion: presentar, explicar, comparar, valorar o
%   concluir.
% - Usar primera persona con moderacion academica: "Considero que",
%   "Desde esta perspectiva", "La posicion que sostengo es". Evitar "yo creo".
%
% Notas de enfoque personal segun enfoque_personal.txt:
% - La voz debe ser academica, tecnica, directa, estrategica y funcional.
% - El texto no debe sonar como estudiante pasivo que recopila: debe sonar como
%   alguien que interpreta, ordena, cruza datos y busca consecuencia practica.
% - Antes de redactar, identificar la funcion de cada seccion: abrir problema,
%   definir criterio, demostrar con evidencia, comparar, valorar o concluir.
% - Evitar frases genericas como "es importante porque si" o "desde siempre".
%   Sustituirlas por relaciones verificables: causa, efecto, limite, costo,
%   funcion, riesgo, implicacion o mejora.
% - Cada concepto debe pasar de definicion a analisis. Formula util:
%   concepto -> funcion -> riesgo si se aplica mal -> consecuencia -> postura.
% - La postura personal no debe ser opinion suelta. Debe tener tres piezas:
%   posicion, razon y efecto. Ejemplo: "Considero que X, porque Y; por ello Z".
% - Integrar experiencia o criterio propio sin perder formalidad: no narrar la
%   vida personal, sino usar mirada practica para explicar por que el tema
%   importa en educacion, derecho, instituciones, comunicacion o profesion.
%
% Estilo observado en actividades ya realizadas:
% - Abrir con una tesis concreta:
%   "El Derecho no solo ordena conductas: organiza poder, define limites y
%   distribuye consecuencias."
% - Formular el problema en negativo y positivo:
%   "El problema no es solo..., sino..."
% - Convertir el producto visual en argumento:
%   "El cuadro no solo recopila datos; reconstruye la funcion, el problema y el
%   punto que requiere correccion."
% - Cerrar con aprendizaje transferible:
%   "No basta con saber que dice la norma; tambien debe analizarse a quien
%   protege, a quien excluye y que razones permiten defenderla publicamente."
%
% Nivel cognitivo esperado:
% - Recordar: usar solo para definir conceptos basicos. No debe ser el nivel
%   dominante en una actividad ponderable.
% - Comprender: explicar con palabras propias y relacionar con el tema.
% - Aplicar: llevar el concepto a un caso, norma, texto, correo, dilema o hecho.
% - Analizar: separar elementos, detectar causas, efectos, tensiones y limites.
% - Evaluar: sostener una postura con criterios, fuentes y consecuencias.
% - Crear: elaborar producto propio: mapa, cuadro, linea de tiempo, ensayo,
%   propuesta, reformulacion, glosario, portafolio o solucion argumentada.
%
% Regla de nivel minimo:
% - Si la actividad vale puntos y pide producto, no quedarse en recordar o
%   comprender. Debe llegar por lo menos a aplicar y analizar; si pide conclusion
%   o postura, debe llegar a evaluar.
%
% Uso de las 100 tecnicas didacticas:
% - Esta plantilla conserva el manual "Tecnicas didacticas del aprendizaje" que
%   se incorporo en las plantillas de Filosofia del Derecho, Etica y Moral
%   juridica, y Redaccion en contextos virtuales.
% - Cuando la planeacion mencione una tecnica de la coleccion UnADM, identificar:
%   definicion de la tecnica, producto esperado, pasos de construccion, criterios
%   de revision y cierre reflexivo.
% - La tecnica didactica no debe verse como adorno. Debe mostrar seleccion,
%   jerarquia, relacion entre ideas y lectura propia.
% - Siempre que sea viable, construir la tecnica en LaTeX/TikZ: cuadro
%   comparativo, mapa conceptual, esquema, mapa mental, linea de tiempo, matriz,
%   red conceptual, diagrama de flujo, tabla SQA o lista de verificacion.
%
% Checklist antes de entregar:
% [ ] El titulo, materia, semana, docente y fecha estan actualizados.
% [ ] El objetivo coincide con la planeacion.
% [ ] El producto solicitado aparece visible, rotulado y completo.
% [ ] La introduccion plantea un problema o postura, no solo anuncia el tema.
% [ ] El desarrollo explica, conecta y analiza.
% [ ] Hay citas APA autor-anio dentro del texto.
% [ ] La bibliografia final coincide con las fuentes citadas.
% [ ] La conclusion responde que se comprendio y que postura se sostiene.
% [ ] El nivel cognitivo llega al menos a aplicacion y analisis.
% [ ] Si hay conclusion/postura, incluye posicion, razon y consecuencia.
% [ ] La tecnica didactica solicitada esta construida y explicada.
% [ ] No aparecen instrucciones copiadas de la planeacion dentro del producto.
% [ ] Ortografia, acentos, sintaxis y nombres propios revisados.
% [ ] Si se uso IA como apoyo, el texto fue comprendido, revisado y apropiado.
% ==============================================================================

% CREACION DEL DOCUMENTO
\documentclass[
	spanish,
	letterpaper, oneside
]{article}

% ==============================================================================
% INFORMACION DEL DOCUMENTO
% Actualizar en cada actividad.
% Ejemplos:
% \def\documenttitle {Cuadro comparativo de conceptos fundamentales}
% \def\documentsubtitle {Actividad 1 - Contabilidad I}
% \def\documentsubject {Actividad 1}
% ==============================================================================
\def\documenttitle {Conceptos basicos de la contabilidad}
\def\documentsubtitle {Cuestionario resuelto: ciclo contable, partida doble y ecuacion contable}
\def\documentsubject {Actividad 1 - Contabilidad I}

\def\documentauthor {Martin Jonathan de la Cruz}
\def\coursename {Contabilidad I}
\def\coursecode {CON-I}

\def\universityname {Universidad Ciudadana de Nuevo Leon}
\def\universityfaculty {Programa academico UCNL}
\def\universitydepartment {Contabilidad I}
\def\universitydepartmentimage {departamentos/logo-ucnl}
\def\universitydepartmentimagecfg {height=1.57cm}
\def\universitylocation {Monterrey, Nuevo Leon}

% INTEGRANTES, PROFESORES Y FECHAS
\def\authortable {
	\begin{tabular}{ll}
		\textbf{Alumno:}
		& \begin{tabular}[t]{l}
			 Martin Jonathan de la Cruz \\
		\end{tabular} \\
		Matricula: & UCNL \\

		\textit{Figura docente:}
		& \begin{tabular}[t]{l}
			Nombre \\ Apellido
		\end{tabular} \\
		& \\
		\multicolumn{2}{l}{Fecha de realizacion: \today} \\
		\multicolumn{2}{l}{\universitylocation}
	\end{tabular}
}

% IMPORTACION DEL TEMPLATE BASE
\input{template}

% FORMATO GENERAL
\usepackage{helvet}
\renewcommand{\familydefault}{\sfdefault}

% TABLAS Y PRODUCTOS VISUALES
\usepackage{array}
\usepackage{longtable}
\usepackage{pdflscape}
\usepackage{tikz}
\usetikzlibrary{
	arrows.meta,
	positioning,
	calc,
	matrix,
	fit,
	shapes.geometric,
	shadows.blur
}

% CITAS EN FORMATO AUTOR-ANIO, COMPATIBLE CON APA 7 EN EL CUERPO DEL TEXTO
% Usar:
% - Cita parentetica: \citep{clave}
% - Cita narrativa: \citet{clave}
% - Varias fuentes: \citep{clave1,clave2}
\setcitestyle{authoryear,open={(},close={)}}

% ==============================================================================
% AYUDAS DE REDACCION
% Quitar o comentar \pendiente cuando el documento final este terminado.
% ==============================================================================
\newcommand{\pendiente}[1]{}

% ==============================================================================
% MARCA DE AGUA EN PORTADA
% - Se aplica solo a la portada porque usa \AddToShipoutPictureBG*.
% - Para apagarla en alguna actividad: \def\coverwatermarkenabled {false}
% - Usar opacidad baja para que no compita con titulo, datos ni logotipo.
% ==============================================================================
\def\coverwatermarkenabled {true}
\def\coverwatermarkimage {img/departamentos/logo-nuevo-leon.png}
\def\coverwatermarkopacity {0.16}
\def\coverwatermarkwidth {0.72\paperwidth}
\def\coverwatermarkxshift {0cm}
\def\coverwatermarkyshift {-0.35cm}

\newcommand{\insertcoverwatermark}{%
	\ifthenelse{\equal{\coverwatermarkenabled}{true}}{%
		\AddToShipoutPictureBG*{%
			\begin{tikzpicture}[remember picture,overlay]
				\node[
					opacity=\coverwatermarkopacity,
					inner sep=0pt,
					xshift=\coverwatermarkxshift,
					yshift=\coverwatermarkyshift
				] at (current page.center) {%
					\includegraphics[width=\coverwatermarkwidth]{\coverwatermarkimage}%
				};
			\end{tikzpicture}%
		}%
	}{}%
}

\begin{document}

% PORTADA
\insertcoverwatermark
\templatePortrait

% CONFIGURACION DE PAGINA Y ENCABEZADOS
\templatePagecfg
\onehalfspacing

% ==============================================================================
% RESUMEN / ABSTRACT
% Debe ser breve. Formula recomendada:
% Tema + objetivo + tecnica/producto + enfoque personal.
% ==============================================================================
\begin{abstractd}
	\textbf{Objetivo:} Identificar conceptos basicos del ciclo contable, la clasificacion de cuentas y la logica de la partida doble mediante la resolucion de un cuestionario de opcion multiple.

	\textbf{Producto elaborado:} Cuestionario resuelto en formato de tabla sobre registros contables, libro mayor, balanza de comprobacion, ecuacion contable y clasificacion de cuentas.

	\textbf{Enfoque de analisis:} El cuestionario se analiza como tecnica diagnostica: cada reactivo permite verificar si el concepto contable se entiende por su funcion dentro del registro, la clasificacion y el control de la informacion financiera.

	\textbf{Nivel cognitivo:} Comprender y aplicar, porque el cuestionario exige reconocer conceptos basicos y seleccionar la respuesta correcta segun su efecto en cuentas, libros y estados financieros.
\end{abstractd}

% TABLA DE CONTENIDOS - INDICE
\templateIndex

% CONFIGURACIONES FINALES
\templateFinalcfg

% ======================= INICIO DEL DOCUMENTO =======================

\section{Introduccion}

La contabilidad no es solamente una lista de cuentas: es un sistema que ordena hechos economicos para convertirlos en informacion util. Por eso, conceptos como cargo, abono, libro mayor, balanza de comprobacion y ecuacion contable deben entenderse por la funcion que cumplen dentro del proceso de registro y verificacion \citep{openstaxFinancialAccounting2019,romeroContabilidad2018}.

La dificultad principal consiste en reconocer que una respuesta contable es correcta cuando conserva la logica del ciclo contable. Una cuenta puede aumentar o disminuir segun su naturaleza; una balanza no reemplaza los estados financieros, sino que verifica equilibrio; y la partida doble sostiene la relacion entre cargos y abonos \citep{romeroContabilidad2018,kiesoIntermediateAccounting2020}.

El tema se ubica en la base conceptual de Contabilidad I. Antes de elaborar registros complejos, es necesario dominar la clasificacion de cuentas, la ecuacion contable, la partida doble y la ruta que siguen las operaciones desde el diario hasta el mayor y la balanza de comprobacion. Si se clasifica mal una cuenta, se altera el saldo; si no se comprende la partida doble, se pierde equilibrio; y si se confunde la balanza con los estados financieros, se interpreta de forma incompleta el cierre del periodo \citep{cinifNif2025,openstaxFinancialAccounting2019}.

% ENCABEZADO TEMATICO (no 'Desarrollo'): nombra el contenido, no el proceso.
\section{Fundamentos del ciclo contable y la partida doble}

El \textbf{ciclo contable} organiza el registro de transacciones, su clasificacion por cuentas, la verificacion de saldos y la preparacion de estados financieros. La \textbf{ecuacion contable}, expresada como Activo $=$ Pasivo + Capital, muestra que los recursos de una entidad se financian con obligaciones o con aportaciones y resultados acumulados. Esta relacion es consistente con la estructura general de la informacion financiera establecida por las Normas de Informacion Financiera mexicanas \citep{cinifNif2025,imcpNif2024}.

La \textbf{partida doble} sostiene que cada \textbf{cargo} debe corresponder con un \textbf{abono} equivalente. En consecuencia, el cargo no siempre aumenta ni siempre disminuye: su efecto depende de la naturaleza de la cuenta. El \textbf{libro diario} conserva el orden cronologico de las operaciones, el \textbf{libro mayor} las organiza por cuenta y la \textbf{balanza de comprobacion} verifica el equilibrio antes del cierre; ninguno reemplaza a los estados financieros. Esta regla explica por que el cuestionario debe resolverse con criterio tecnico y no con respuestas mecanicas \citep{romeroContabilidad2018,kiesoIntermediateAccounting2020}.

% Nivel cognitivo aplicado para el motor-inteligente:
% Comprender y aplicar. Comprender implica reconocer que representa cada cuenta o
% documento; aplicar exige seleccionar la respuesta correcta segun el efecto que
% produce dentro del sistema contable. La justificacion evita respuestas mecanicas.

Cada reactivo se lee a partir de tres elementos: pregunta, respuesta correcta y razon tecnica.

\begingroup
\scriptsize
\setlength{\tabcolsep}{3.2pt}
\renewcommand{\arraystretch}{1.28}
\begin{longtable}{@{}>{\centering\arraybackslash}p{0.055\textwidth}>{\raggedright\arraybackslash}p{0.31\textwidth}>{\raggedright\arraybackslash}p{0.23\textwidth}>{\raggedright\arraybackslash}p{0.34\textwidth}@{}}
\caption{Cuestionario resuelto de conceptos basicos de Contabilidad I}\label{tab:cuestionario-contabilidad-1}\\
\toprule
\textbf{No.} & \textbf{Pregunta} & \textbf{Respuesta correcta} & \textbf{Justificacion tecnica} \\
\midrule
\endfirsthead
\toprule
\textbf{No.} & \textbf{Pregunta} & \textbf{Respuesta correcta} & \textbf{Justificacion tecnica} \\
\midrule
\endhead
\bottomrule
\endfoot
\bottomrule
\endlastfoot
\midrule
1 & \textquestiondown Que ocurre si una cuenta se carga? & Depende del tipo de cuenta. & El cargo aumenta algunas cuentas y disminuye otras; su efecto depende de la naturaleza contable de la cuenta \citep{romeroContabilidad2018}. \\
\midrule
2 & \textquestiondown En que etapa se trasladan los saldos al libro mayor? & Despues del registro en el diario. & Primero se registra cronologicamente la operacion y despues se concentra por cuenta en el mayor \citep{openstaxFinancialAccounting2019}. \\
\midrule
3 & \textquestiondown Cual es el objetivo principal de la balanza de comprobacion? & Verificar el equilibrio contable. & La balanza comprueba que la suma de cargos y abonos mantenga igualdad antes de preparar informacion final \citep{kiesoIntermediateAccounting2020}. \\
\midrule
4 & \textquestiondown Que establece el principio de partida doble? & Cada cargo debe tener un abono equivalente. & Toda transaccion afecta al menos dos cuentas y conserva equilibrio entre origen y aplicacion de recursos \citep{romeroContabilidad2018}. \\
\midrule
5 & \textquestiondown Que contiene cada cuenta en el libro mayor? & Cargos, abonos y saldo. & El mayor acumula movimientos por cuenta y permite conocer el saldo resultante. \\
\midrule
6 & \textquestiondown Que cuentas se incluyen en el activo? & Recursos y bienes de la empresa. & El activo agrupa recursos controlados por la entidad que pueden generar beneficios economicos. \\
\midrule
7 & \textquestiondown Cual es la etapa final del ciclo contable? & Cierre de cuentas y elaboracion de estados financieros. & El cierre integra resultados y permite preparar informacion financiera del periodo. \\
\midrule
8 & \textquestiondown Cual es la ecuacion contable basica? & Activo $=$ Pasivo + Capital. & La ecuacion muestra que los recursos se financian con obligaciones o patrimonio. \\
\midrule
9 & \textquestiondown Que tipo de cuenta es ``Equipos de oficina''? & Activo fijo. & Representa un bien de uso durable destinado a la operacion de la entidad. \\
\midrule
10 & \textquestiondown Cual es una cuenta de pasivo? & Documentos por pagar. & Los documentos por pagar representan obligaciones de pago frente a terceros. \\
\end{longtable}
\endgroup

% Criterio de organizacion para el motor-inteligente:
% La tabla diagnostica conceptos clave, ordena respuestas y asocia cada opcion
% correcta con una justificacion observable; no debe mostrarse como seccion
% visible salvo que la consigna solicite explicar la tecnica didactica.


% CONCLUSION en pagina propia (\clearpage); analisis+postura integrados; IA como footnote.
\clearpage
\section{Conclusion}

El cuestionario permite reconocer los fundamentos del ciclo contable, la partida doble, la clasificacion de cuentas y la ecuacion contable. La lectura conjunta muestra que los conceptos basicos de contabilidad operan como sistema: la partida doble explica la relacion entre cargos y abonos; el libro diario conserva el orden cronologico; el mayor organiza la informacion por cuenta; y la balanza de comprobacion revisa que exista equilibrio antes del cierre. La solidez de cada respuesta se evaluo con criterios claros: su correspondencia con la logica del ciclo contable y su fundamento en la normativa consultada \citep{openstaxFinancialAccounting2019,kiesoIntermediateAccounting2020}, no la mera memoria.

Desde mi analisis, la postura que sostengo es que Contabilidad I debe aprenderse como lenguaje de control y decision, no como memorizacion de cuentas. La razon es que cada registro produce una consecuencia sobre la lectura del patrimonio, las deudas y los resultados; por ello, considero que justificar cada respuesta fortalece la capacidad de interpretar informacion financiera con orden, precision y responsabilidad \citep{cinifNif2025,imcpNif2024}.\footnote{Declaracion de uso de inteligencia artificial: se emplearon herramientas de IA como apoyo para organizar y depurar la redaccion; el analisis conceptual, la seleccion de respuestas y la postura personal son propios, revisados y validados por el autor, sin que la IA sustituyera el criterio academico.}

\clearpage
\bibliography{contabilidad-I}

\end{document}


